\section{Local determination and exterior scattering data}
\label{sec:v71-local-observation}

The endpoint law is constructed from actual finite trajectories.  We now
make precise its dependence on a neighborhood of the marked itinerary.
This also fixes the comparison with the global scattering observation
in Section~\ref{sec:v71-lens-comparison}.

\subsection{The local observation record}

Fix the polygon, its marked frames, a sufficiently small common graph
interval $I$, and the two offsets of
Theorem~\ref{thm:v64-physical-law}.  Shrink $I$ once, if necessary, so
that the stationary bridges for all returning $N\ge N_0$ lie in prescribed
contact collars with compact closure.  There are finitely many such
collars modulo the lattice.  Their successive chords, together with the
short initial and terminal free segments allowed by the time-split
construction, lie in a fixed neighborhood $U$ of the itinerary modulo
the lattice.  We require the same clearance and free-segment margins as
in that theorem.  The neighborhood includes the reflecting collars;
away from them its relevant trajectory tubes are obstacle-free.

A \emph{local completion} consists of closed obstacles agreeing with
the prescribed contact graphs on these collars and satisfying these
margins on $U$; the boundary away from $U$ may vary.  The lattice and marks
are fixed.  This definition concerns actual tables, not independently
assigned Taylor coefficients.  For a chosen phase $b$, returning length
$N$ and offset $d$, a reset preparation is drawn from the normalized
Liouville measure of that table.  Record the endpoint mark $(u,v)$ on
success of the stipulated itinerary event, and a single symbol $\dagger$
otherwise.  Write $\mu_{N,b,d}$ for the success subprobability measure on
the gate $Q_+$, $p_{N,b,d}=\mu_{N,b,d}(Q_+)$, and
\begin{equation}\label{eq:v71-record}
 \mathsf P_{N,b,d}
   =(1-p_{N,b,d})\delta_\dagger+\mu_{N,b,d}.
\end{equation}
The conditional endpoint density is that of $\mu_{N,b,d}/p_{N,b,d}$.
The symbol $\dagger$ contains no trajectory trace or additional data from
failed preparations.  Both the resetting and this recording convention
are essential to the statement below.

\begin{proposition}[Exact locality of the charged record]
\label{prop:v71-locality}
Let $\Omega$ and $\widetilde\Omega$ be two local completions with the
same contact graphs on the prescribed collars, and free areas
$\mathcal A$ and $\widetilde{\mathcal A}$ in the common cell.  Then, for
every phase, both selected offsets, and every returning $N\ge N_0$,
\begin{equation}\label{eq:v71-area-scaling}
 \widetilde\mu_{N,b,d}
       =\frac{\mathcal A}{\widetilde{\mathcal A}}\mu_{N,b,d}.
\end{equation}
In particular the conditional finite-flight densities and their limiting
laws coincide.  If $\mathcal A=\widetilde{\mathcal A}$, the complete
records \eqref{eq:v71-record} coincide, and so do all finite products of
these independently reset experiments.  The assertion is exact, not an
asymptotic estimate as $N\to\infty$.
\end{proposition}
\begin{proof}
The single-edge length functions between successive collars are the
same in the two tables.  The stationary equations and their boundary
conditions are consequently identical.  The weighted uniqueness in
Theorem~\ref{thm:v64-periodic-relative}, on the chosen common collars, identifies
their finite stationary bridges.  They remain physical in each table by
the common clearance margins.  Thus their stationary lengths $W_{N,b}$,
gauged actions $E_{N,b}$, mixed twists and reference normalizations
$D_{N,b}$ agree as functions of the endpoints.  Equality holds throughout
the endpoint gate, not just to finite order at its origin.

Put $\beta_{N,b}=(-\partial_{uv}W_{N,b})/D_{N,b}$ and
$r_{N,b,d}=d-E_{N,b}+p_bu-p_bv$.  The exact time-split computation of
Theorem~\ref{thm:v64-physical-law} gives
\begin{equation}\label{eq:v71-local-measure}
 \frac{d\mu_{N,b,d}}{du\,dv}
   =\frac{D_{N,b}}{2\pi\mathcal A}
       \beta_{N,b}(u,v)\{r_{N,b,d}(u,v)\}_+.
\end{equation}
Indeed the allowed initial free time ranges from $0$ to $r_{N,b,d}$;
the free-segment margins ensure that neither an earlier nor an additional
reflection truncates this interval.  Every successful trajectory is
accounted for by this parametrization because the event specifies all
intermediate collisions in the collars and the last-collision rule.
It does not include other itineraries elsewhere in the table.  All
factors in \eqref{eq:v71-local-measure} except the total free area are
common to the two completions.  This proves
\eqref{eq:v71-area-scaling}.  Integrating gives the same factor for
$p_{N,b,d}$, so normalization cancels it.  The positive success floor
ensures that the conditional densities are defined.  Their limits agree
by the relative law.  At equal area the subprobabilities and hence the
failure masses agree; independence gives the final product assertion.
\end{proof}

\begin{remark}
If settings are selected from the permitted experiments by the same
measurable rule based on earlier recorded outcomes and independent
auxiliary randomization, equality also holds for every fixed finite
number of preparations.  Conditional on the history the next setting
has the same distribution and its record kernel is the same; induction
proves the assertion.  This does not extend it to an unreset trajectory,
to recording a rejected ray's path, or to a table-dependent error kernel.
A common recording kernel applied to the exact records preserves equality.
\end{remark}

\subsection{Complete contact information without remote completion}

Two marked tables will be called \emph{contact equivalent} when their
graph functions agree on some common neighborhood of every contact.
Use a germ of the local law experiment: endpoints may be restricted to a
smaller common positive-residual gate, with its normalization specified
anew.  This convention is necessary when only contact germs, rather than
fixed-radius graphs, are given.  In particular, restricting a density
without renormalizing its original larger gate is a different operation.

\begin{corollary}[Local completion and the determined object]
\label{cor:v71-local-completion}
On the marked smooth class of Theorem~\ref{thm:v70-main}, equality of
the two-offset limiting law germs is equivalent to contact equivalence.
Here equality of laws is tested on a sufficiently small common gate at
every phase.  In the equal-area subclass, contact equivalence also gives
equality of all charged finite-flight records on sufficiently small
common gates.

There are nonconstant smooth one-parameter families of closed strictly
convex obstacles in a fixed periodic cell, with the same free area and
marked polygon, for which these charged records are identical although
the obstacle boundaries away from the collars differ.  These families
can be chosen so that their finite compact obstacle unions in the plane
have different almost-everywhere exterior travelling-time data and
different almost-everywhere scattering length spectra.
\end{corollary}
\begin{proof}
Equality of the limiting laws gives equality of the anchored action
functions by Proposition~\ref{prop:v65-two-offset}.  The global
quadratic inverse, signed finite-jet inversion, and actual smooth
inverse of Theorem~\ref{thm:v68-smooth-rigidity} then identify the graph
functions on a common smaller collar.  These steps use actual realizable
law tuples.  Conversely, if the contact germs agree, choose compact
collars inside their common domains.  The weighted finite and half-line
constructions permit a sufficiently small common endpoint gate for which
all stationary visits stay inside these collars.  Both tables have
positive clearance along the fixed polygon, so the short endpoint
segments also have common margins.  The two offsets are chosen below
these common free-flight bounds, as required by the experiment.  Proposition~\ref{prop:v71-locality} applies.  It proves the
limiting assertion, and the finite-record assertion when the areas agree.
This is an equivalence of germs of experiments; it asserts no equality
on a larger gate outside the common graph domains.

For an actual family take the three-obstacle realization of
Proposition~\ref{prop:v64-triangle}.  Choose $U$ sufficiently small around
its contact collars and flight tubes.  On an obstacle boundary choose
two disjoint open arcs with closures outside $U$.  Such arcs exist on
the side facing away from the interior contact polygon in this explicit
configuration.  Let $\psi_1,\psi_2$ be nonzero smooth nonnegative normal
bumps supported in these arcs.  For small real parameters $(s,t)$ deform
the boundary normally by $s\psi_1+t\psi_2$, leaving all other boundary
points unchanged.  Disjoint compact embedded boundaries stay disjoint
and strictly convex under sufficiently small $C^2$ perturbations; the
lattice translates remain separated as well.  The positive distance
between the bump supports and the closed trajectory tubes preserves
all the required local margins.

Let $V(s,t)$ be the total obstacle area in the cell.  Normal variation of
area gives
\[
 \partial_t V(0,0)=\int\psi_2\,d\sigma>0.
\]
For completeness this derivative follows by using boundary normal
coordinates: a strip of signed width $t\psi_2(\sigma)$ has area
$t\int\psi_2\,d\sigma+O(t^2)$, with a uniform remainder on the compact
support.  The area is a smooth function of the two parameters.  The
implicit function theorem gives a smooth $t=t(s)$ near zero, with
$t(0)=0$ and $V(s,t(s))=V(0,0)$.  The family is nonconstant, since on
the first support its normal displacement is exactly $s\psi_1$ and
the second bump has disjoint support.  All contact graphs and marks are
unchanged.  The free area is constant, so
Proposition~\ref{prop:v71-locality} proves equality of every indicated
finite record.  The observation samples are still drawn from the full
table, not from a law postulated only on $U$.

Regard the same three obstacles, without their lattice copies, as a
compact union $K_s$ in $\R^2$.  All $K_s$ lie in one fixed enclosing
circle and satisfy the strictly convex $C^3$ hypotheses of
Noakes--Stoyanov \cite[Theorem~1.1]{NoakesStoyanov2020}.  For $s\ne s'$
sufficiently small, $K_s\ne K_{s'}$ by the disjoint-support construction.
That theorem implies that their travelling-time families are not equal
on a full-measure set of incoming positions and directions.  Its
separate scattering-length conclusion gives the analogous distinction
for incoming/outgoing direction pairs.  The periodic records remain
identical by the preceding paragraph.  This proves the asserted
comparison on an explicitly common class of geometric realizations.
\end{proof}

The two parts of this comparison serve different purposes.  In the local
experiment, the inverse determines the actual visited contact germs,
including flat changes, whereas the construction varies only unvisited
boundary pieces and leaves even the finite charged record fixed.  The
exterior observation distinguishes those pieces by a different data
family.  Hence equality of the local record cannot simply be substituted
for the equality hypothesis in a global lens theorem.  There is no claim
that local laws distinguish equal full marked length spectra, nor that a
known whole obstacle and the supplied marks could not generate its local
laws.  The comparison sharpens the local determination statement without
weakening it and without identifying the two observation models.
